\documentclass[11pt,a4paper]{article}
\usepackage[T1]{fontenc}
\usepackage[utf8]{inputenc}
\usepackage[polish]{babel}
\usepackage{lmodern}
\usepackage[a4paper,margin=26mm]{geometry}
\usepackage[hidelinks]{hyperref}
\title{Postrzegane i obiektywne bezpieczeństwo kontaktu człowieka z humanoidem}
\author{}
\date{Prospektywny szkic programu arc-robo -- MSC-HRI-P08}
\begin{document}
\maketitle
\begin{abstract}
Dokument stanowi lokalny plan badawczy, a nie kopię manuskryptu źródłowego i nie przedstawia wyników empirycznych. Przedmiot: Postrzegane i obiektywne bezpieczeństwo kontaktu człowieka z humanoidem.
\end{abstract}
\section{Pytanie badawcze i cel}
Ocenić sposób weryfikacji problemu: Postrzegane i obiektywne bezpieczeństwo kontaktu człowieka z humanoidem. Hipoteza kierunkowa, operacjonalizacja i kryteria decyzji wymagają zatwierdzenia w protokole repozytorium źródłowego.
\section{Wynik główny i granice}
Wskaźnik pierwszorzędowy oraz punkt odniesienia wymagają określenia w źródłowym protokole.

Granice wnioskowania i niezależność tego manuskryptu wymagają weryfikacji względem prac sąsiednich.
\section{Plan materiału i metod}
Walidacja offline na zamrożonych scenariuszach interakcji humanoidalnej, przy porównywalnych warunkach i niezależnym systemie bezpieczeństwa; badania fizyczne wyłącznie po osobnej kwalifikacji. Należy określić jednostkę analizy, punkty odniesienia, niepewność i kryteria wykluczeń przed pomiarem.
\section{Plan analizy}
Przed analizą należy zdefiniować metryki, podziały niezależne od jednostki eksperymentalnej, przedziały niepewności oraz analizę wrażliwości. Nie przenosić wyników symulacji na certyfikację bezpieczeństwa sprzętu.
\section{Wyniki}
Brak wyników: eksperyment nie został zweryfikowany na potrzeby tego szkicu. Nie wpisywać statystyk bez dowodów źródłowych.
\section{Ograniczenia i etyka}
Dostępność aparatury, licencje, kwalifikacja bezpieczeństwa i ewentualna zgoda na badania osób wymagają odrębnych dowodów.
\section{Proweniencja}
Kanoniczny manuskrypt/protokół pozostaje w projekcie źródłowym. Repozytorium źródłowe ma ograniczony dostęp; publiczny szkic nie zawiera adresu prywatnego źródła. Karta \texttt{study.md} opisuje relacje stanowisk, aplikacji i bibliotek.
\end{document}
